\documentclass[../main.tex]{subfiles}

\begin{document}

\subsection{Overview}

CVE-2024-58100 is a \textbf{state-tracking vulnerability} in the Linux eBPF verifier related to the \texttt{changes\_pkt\_data} property of BPF subprograms. The Linux CNA identifies affected released kernels from 5.6 through 6.6.89 and from 6.7 through 6.12.24; Linux 6.6.90, 6.12.25 and 6.13 contain the corresponding fixes. The bug appears when an \texttt{EXT} program replaces a global BPF subprogram without a proper compatibility check on packet-data side effects.

The original caller may be verified under the assumption that the called global subprogram does not modify packet data. If the replacement \texttt{EXT} program instead calls a helper such as \texttt{bpf\_skb\_change\_head()} or \texttt{bpf\_skb\_pull\_data()}, the underlying \texttt{skb} buffer may be reallocated while the caller still reuses a previously validated \texttt{PTR\_TO\_PACKET}.

The vulnerable pattern is conceptually simple. A traffic-control BPF program stores \texttt{ctx->data} into a register, usually \texttt{r6}, and validates it against \texttt{ctx->data\_end}. The verifier then marks \texttt{r6} as a valid packet pointer. The program subsequently calls a global subprogram that invokes a packet-changing helper such as \texttt{bpf\_skb\_change\_head()}. At runtime, this helper frees and reallocates \texttt{skb->head}. However, because the verifier does not properly reject an \texttt{EXT} replacement with incompatible \texttt{changes\_pkt\_data} behavior, the caller can still be verified under the stale assumption that packet data remains unchanged across the call.

As a result, the BPF program can read and write through a stale packet pointer. In the packet layout used by the local experiment, runtime execution operates on a freed \texttt{kmalloc-1024} slab object, producing a \textbf{Use-After-Free read/write primitive}.

\begin{table}[H]
\centering
\begin{tabularx}{\textwidth}{|l|X|}
\hline
\textbf{CVE ID} & CVE-2024-58100 \\
\hline
\textbf{Affected versions} & Linux 5.6--6.6.89 and 6.7--6.12.24 \\
\hline
\textbf{Vulnerability type} & Verifier state-tracking bug, stale \texttt{PTR\_TO\_PACKET}, UAF \\
\hline
\textbf{Root cause} & Missing compatibility check between an EXT program and the \texttt{changes\_pkt\_data} property of the global subprogram it replaces \\
\hline
\textbf{Target slab} & \texttt{kmalloc-1024} \\
\hline
\textbf{Impact} & Official bug can expose stale packet pointers; full LPE demonstrated from the equivalent local UAF primitive in the controlled environment \\
\hline
\end{tabularx}
\caption{Summary of CVE-2024-58100.}
\end{table}

The post-UAF exploitation chain was evaluated on Linux 6.12.24. The local programs load one \texttt{SCHED\_CLS} object containing a global subprogram. They do not attach a separate EXT replacement and must therefore be treated as \textbf{surrogate reproducers of the stale-pointer primitive}, not as complete demonstrations of the official EXT/\texttt{freplace} sequence. Starting from that primitive, the first exploit reliably recovered its UAF sentinel but depended on a sensitive cross-CPU reclaim race. The second exploit replaces that race with SLUB freelist poisoning and repeatedly completes the privilege-escalation chain in the controlled VM.

\subsection{Vulnerability Analysis}

In traffic-control programs, packet access normally follows a strict verifier pattern:

\begin{enumerate}
    \item read \texttt{ctx->data};
    \item read \texttt{ctx->data\_end};
    \item check that \texttt{data + len <= data\_end};
    \item access the packet through the validated pointer.
\end{enumerate}

The verifier then treats the register holding \texttt{ctx->data} as a bounded \texttt{PTR\_TO\_PACKET}. If a helper later modifies the underlying packet buffer, all old packet pointers must be invalidated. This is normally handled through the \texttt{changes\_pkt\_data} flag.

The official bug appears when an EXT program replaces a global subprogram with incompatible packet-data side effects. Conceptually, the original target has no packet-changing side effect, while its replacement does:

\begin{lstlisting}[language=C, caption={Vulnerable high-level pattern.}]
SEC("classifier")
int main_cls(struct __sk_buff *skb)
{
    void *data     = (void *)(long)skb->data;
    void *data_end = (void *)(long)skb->data_end;

    if (data + 128 > data_end)
        return TC_ACT_OK;

    /* data is now verifier-approved */
    target_subprog(skb); /* replaced at runtime through freplace */

    /*
     * Verifier still trusts data.
     * Runtime: data points to freed skb->head.
     */
    *(u32 *)(data + 0) = 0x41414141;

    return TC_ACT_OK;
}

__noinline int target_subprog(struct __sk_buff *skb)
{
    return 0; /* contract used while verifying the caller */
}

/* Separate BPF_PROG_TYPE_EXT replacement */
int replacement(struct __sk_buff *skb)
{
    return bpf_skb_change_head(skb, 1024, 0);
}
\end{lstlisting}

The key issue is not only the local behavior of the packet-changing helper, but the missing compatibility check between the original global subprogram and the \texttt{EXT} program that replaces it. The caller is verified under one packet-data contract, while runtime execution may follow a replacement implementation with different side effects. Runtime behavior is different: \texttt{bpf\_skb\_change\_head()} frees the old \texttt{skb->head}, so the caller's packet pointer becomes dangling.

The vulnerability is therefore another verifier/runtime divergence:

\begin{itemize}
    \item \textbf{Verifier view}: \texttt{r6} is still a valid bounded packet pointer.
    \item \textbf{Runtime view}: \texttt{r6} points into a freed \texttt{kmalloc-1024} object.
\end{itemize}


\subsection{Exploit Environment}
CVE-2024-58100 required a specific boot-time setup because the exploit transforms a stale packet pointer into corruption of the SLUB allocator and, finally, into control of a \texttt{pipe\_buffer}. The target was Linux 6.12.24 with \texttt{CONFIG\_KASAN=n} and \texttt{CONFIG\_SLAB\_FREELIST\_HARDENED=n}.

The exploit also needs access to BPF program loading, network-related BPF functionality, kernel symbol addresses and a writable directory for the final \texttt{modprobe\_path} payload.
For this reason, the Buildroot image includes an init script named \texttt{S99exploit}. This script runs automatically at VM boot and prepares the environment before the exploit is executed.
The script performs the following operations:
\begin{itemize}
    \item ensures that \texttt{/tmp} exists and is world-writable with mode \texttt{1777};
    \item sets \texttt{/proc/sys/kernel/kptr\_restrict} to \texttt{0}, allowing real kernel addresses to be read from \texttt{/proc/kallsyms};
    \item sets \texttt{/proc/sys/kernel/unprivileged\_bpf\_disabled} to \texttt{0} as part of the permissive laboratory setup;
    \item lowers \texttt{/proc/sys/kernel/perf\_event\_paranoid}; this is part of the general test setup, while symbol visibility is provided by \texttt{kptr\_restrict=0} together with the assigned capabilities;
    \item copies the PoC and the two exploit variants, including \texttt{exploit2}, into \texttt{/home/user};
    \item assigns the required capabilities to the copied binaries using \texttt{setcap}.
\end{itemize}
The relevant capabilities assigned to the user binaries are:
\begin{lstlisting}[language=bash, caption={Capabilities assigned by S99exploit.}]
cap_bpf,cap_net_admin,cap_perfmon,cap_syslog+ep
\end{lstlisting}
This setup allows a non-root user to launch the exploit while carrying privileged capabilities. It is therefore a \textbf{non-root but capability-bearing} execution context, not an unprivileged one. In particular, \texttt{cap\_bpf} enables privileged BPF loading, \texttt{cap\_net\_admin} covers network-related BPF operations, \texttt{cap\_syslog} permits access to restricted kernel-address information, and \texttt{cap\_perfmon} is included by the setup for BPF observability and compatibility. The process is pinned to CPU 0 during heap manipulation so that allocations use a consistent per-CPU SLUB freelist.

The addresses of \texttt{modprobe\_path} and \texttt{anon\_pipe\_buf\_ops} are read from \texttt{/proc/kallsyms}. This is a laboratory convenience enabled by \texttt{kptr\_restrict=0}, not a KASLR bypass derived from the vulnerability. The complete configuration is intentionally permissive and should not be interpreted as a hardened production setup.

\subsection{Exploit Primitive}

The local surrogate reproducer provides a \textbf{stable UAF read/write on a freed \texttt{kmalloc-1024} slot}. This slab class is relevant because it is also used by arrays of \texttt{struct pipe\_buffer}. A default pipe uses \texttt{PIPE\_DEF\_BUFS = 16} entries, and each \texttt{struct pipe\_buffer} is 40 bytes. Therefore:

\begin{lstlisting}[caption={\texttt{pipe\_buffer} allocation size.}]
16 * sizeof(struct pipe_buffer)
= 16 * 40
= 640 bytes
-> rounded by SLUB to kmalloc-1024
\end{lstlisting}

Both exploit variants start from this same UAF primitive. They differ in how they try to turn control over the freed slot into corruption of a live kernel object.

\subsection{Original Race-Based Exploit}

The first exploit tried to reclaim the freed \texttt{skb->head} directly with a \texttt{pipe\_buffer[]} allocation. After calling \texttt{bpf\_skb\_change\_head()}, the BPF program kept the stale packet pointer in \texttt{r6}, executed a busy-wait loop and then wrote the address of a userspace \texttt{fake\_ops} table at \texttt{r6+32}. If another CPU allocated and initialized a pipe-buffer array in the freed slot during that interval, this write targeted \texttt{pipe\_buffer[2].ops}.

Userspace prepared the \texttt{kmalloc-1024} cache with \texttt{msg\_msg} allocations and launched separate BPF and pipe-worker threads on different CPUs. The old source also contains a KASAN-quarantine drain routine for experiments on instrumented kernels; however, KASAN was disabled in the Linux 6.12.24 configuration documented here, so no KASAN quarantine was active during these tests. The pipe workers created pipes and filled three entries so that \texttt{pipe\_buffer[2]} became active. The BPF workers repeatedly freed packet buffers and performed the delayed stale write, while userspace triggered the \texttt{modprobe\_path} payload to detect a successful corruption.

This design reliably confirmed the UAF and its read/write capability, but the privilege-escalation stage depended on a narrow ordering: \texttt{kcalloc()} had to reclaim and zero the slot before the final BPF write, and the write had to occur before the pipe was released. Per-CPU freelists, cross-CPU allocation timing and the small number of virtual CPUs made this race unstable in the test VM. The root fallback validated the final \texttt{modprobe\_path} payload independently, but did not demonstrate a complete non-root escalation through the race.

\subsection{Strategic and Architectural Differences}

The new exploit preserves the stale-pointer trigger and the final \texttt{modprobe\_path} objective, but changes the intermediate goal. The old strategy tried to make two concurrent activities meet on the same freed object at the correct instant. The new strategy first corrupts SLUB metadata and then lets ordinary allocator behavior deliver a future kernel object to an address already controlled by userspace.

\begin{table}[H]
\centering
\begin{tabularx}{\textwidth}{|l|X|X|}
\hline
\textbf{Aspect} & \textbf{Original exploit} & \textbf{New exploit} \\
\hline
Allocator strategy & Direct reuse of the freed slot through a cross-CPU allocation race & Poison the freepointer and redirect a later allocation \\
\hline
Execution model & Parallel BPF and pipe workers with a busy-wait timing window & Single-CPU grooming followed by bounded pipe spraying \\
\hline
Stale write target & \texttt{r6+32}, interpreted as \texttt{pipe\_buffer[2].ops} after successful reuse & \texttt{r6+448}, the SLUB freepointer inside the freed object \\
\hline
Controlled object & A pipe-buffer array that must occupy the UAF slot before the delayed write & A pipe-buffer array allocated directly in the mapped userspace page \\
\hline
Final corruption & BPF writes \texttt{fake\_ops} into the reclaimed kernel slot & Userspace edits \texttt{pipe\_buffer[0].ops} in its own mapped page \\
\hline
Observed result & Stable UAF; complete LPE race not reliable & Repeatable full LPE in the documented VM configuration \\
\hline
\end{tabularx}
\caption{Strategic and architectural differences between the two CVE-2024-58100 exploits.}
\end{table}

Architecturally, the main improvement is the separation of responsibilities. In the original version, the BPF program had to free the object and perform the final operations-pointer overwrite at precisely the right time. In the new version, BPF is used only to stamp the freed object and poison its freepointer. Once the allocation reaches the mapped page, userspace can inspect and modify the resulting \texttt{pipe\_buffer} without a second timing-dependent stale write.

\subsection{SLUB Freelist-Poisoning Strategy}

\textbf{SLUB} is the kernel allocator used to manage caches of similarly sized objects. Each cache maintains a \textbf{freelist}, namely a chain of objects that are currently free and can satisfy future allocations. A free object contains a \textbf{freepointer} that identifies the next object in this chain. \textbf{Freelist poisoning} consists of overwriting that pointer so that a later allocation follows an attacker-selected address.

For the tested \texttt{kmalloc-1024} cache, SLUB stores the freepointer halfway through the object, at offset 512. The stale packet register does not point to the beginning of the allocation: on the tested x86-64 configuration, \texttt{r6} points 64 bytes into the old \texttt{skb->head}, after the packet headroom. The required write offset is therefore:

\begin{lstlisting}[caption={Freepointer offset relative to the stale packet pointer.}]
SLUB freepointer offset = 512
stale r6 offset         = 64
FP_R6_OFF               = 512 - 64 = 448
VERIFY_BYTES             = 448 + 8 = 456
\end{lstlisting}

The BPF payload writes the fixed userspace address \texttt{0x13370000} at \texttt{r6+448}. When SLUB later removes the corrupted object from the freelist, it loads this value as the address of the next free object. A subsequent \texttt{kmalloc-1024} request can therefore return the mapped userspace page as if it were a legitimate slab object.

This provides controlled allocation placement rather than direct object reuse. The technique is still dependent on the tested allocator layout and uses repeated grooming attempts, because the poisoned object may not immediately become the next allocation served by the active per-CPU freelist.

\subsection{Exploit Architecture}

The exploit consists of a userspace orchestrator and a \texttt{SCHED\_CLS} BPF payload containing a global subprogram. The official root cause concerns compatibility between an \texttt{EXT}/\texttt{freplace} program and the global subprogram it replaces. The local code does not implement that attachment boundary: it directly constructs the stale-pointer behavior inside one BPF object so that the UAF and allocator manipulation can be studied. Consequently, the code validates the post-UAF exploitation strategy, while the official EXT compatibility defect is established by source and patch analysis rather than reproduced end to end by this program.

The BPF payload performs the following operations:

\begin{lstlisting}[caption={Simplified BPF payload used by the new exploit.}]
1. Save ctx->data in r6 and ctx->data_end in r7
2. Verify access through r6 + 456
3. Call pkt_mutator(), which invokes bpf_skb_change_head()
4. The old kmalloc-1024 skb->head is freed
5. Write 0x4242424241414141 through stale r6
6. Write 0x13370000 at stale r6 + 448
7. Copy 128 bytes from the stale object into a BPF map
\end{lstlisting}

The sentinel confirms that the stale packet pointer still provides read/write access after the reallocation of \texttt{skb->head}. The second write corrupts the SLUB freepointer.

The complete userspace-assisted exploitation path is summarized below:

\begin{figure}[H]
\centering
\begin{tikzpicture}[
    node distance=0.55cm,
    every node/.style={
        draw,
        rounded corners,
        align=center,
        font=\scriptsize,
        minimum height=0.8cm,
        text width=2.25cm
    },
    arrow/.style={-{Latex[length=2mm]}, thick}
]

\node (setup) {Map target page\\and resolve symbols};
\node[right=of setup] (uaf) {Trigger and\\confirm UAF};
\node[right=of uaf] (poison) {Poison SLUB\\freepointer};
\node[right=of poison] (spray) {Drain freelist\\and spray pipes};

\node[below=0.75cm of spray] (land) {\texttt{pipe\_buffer}\\lands in mmap};
\node[left=of land] (ops) {Replace \texttt{ops}\\with \texttt{fake\_ops}};
\node[left=of ops] (modprobe) {Overwrite\\\texttt{modprobe\_path}};
\node[left=of modprobe] (root) {Create SUID\\root shell};

\draw[arrow] (setup) -- (uaf);
\draw[arrow] (uaf) -- (poison);
\draw[arrow] (poison) -- (spray);
\draw[arrow] (spray) -- (land);
\draw[arrow] (land) -- (ops);
\draw[arrow] (ops) -- (modprobe);
\draw[arrow] (modprobe) -- (root);

\end{tikzpicture}
\caption{CVE-2024-58100 exploitation path using SLUB freelist poisoning.}
\label{fig:path-2024-58100}
\end{figure}

The spray is repeated for a bounded number of iterations because the poisoned object may not immediately be at the head of the active per-CPU freelist. Pipes are used for both draining and reclaiming: if one receives the userspace page, it is deliberately not closed before the corrupted state is neutralized, avoiding an invalid \texttt{kfree()} of userspace memory.

The fake release callback writes the attacker-controlled helper path into \texttt{modprobe\_path}. It also sets \texttt{pipe->ring\_size} to zero and \texttt{pipe->bufs} to \texttt{NULL}; this prevents the normal pipe cleanup path from iterating over or freeing the forged userspace array after control returns to the kernel.

\subsection{Experimental Results}

The new exploit successfully demonstrated the locally reproduced stale-pointer primitive and the complete escalation chain built on top of it. Figure~\ref{fig:log-2024-58100-pe} shows a successful execution from the non-root, capability-bearing test account.

\begin{figure}[H]
    \centering
    \includegraphics[width=0.62\textwidth]{img/CVE-2024-58100-PE.png}
    \caption{Successful CVE-2024-58100 exploitation through SLUB freelist poisoning.}
    \label{fig:log-2024-58100-pe}
\end{figure}

The first part of the log confirms that \texttt{bpf\_skb\_change\_head()} enlarged the packet, freed the old \texttt{kmalloc-1024} head and left \texttt{r6} dangling in the local reproducer. Reading back the sentinel \texttt{0x4242424241414141} then proves that the accepted program wrote through this stale pointer. The log's ``KASLR bypass'' label refers only to reading symbols from the deliberately exposed \texttt{/proc/kallsyms}; it is not a vulnerability-derived KASLR bypass.

In the same run, the allocator followed the poisoned freepointer during the first grooming iteration and placed a pipe-buffer array in the page mapped at \texttt{0x13370000}. The logged value of \texttt{pipe\_buffer[0].ops} matched \texttt{anon\_pipe\_buf\_ops}; userspace then replaced it with \texttt{fake\_ops} and consumed the pipe data to transfer control to the forged release callback.

The resulting \texttt{fake\_release()} invocation changed \texttt{modprobe\_path} in kernel context. Executing the invalid-format trigger then ran the attacker-controlled helper as root. The helper copied the exploit binary, changed its owner to UID 0 and set mode \texttt{4755}. Executing this copy produced a root shell, completing the transition from the non-root test account to effective UID 0.

The exploit was tested repeatedly in the controlled VM and completed the LPE without the old cross-CPU allocation race. This result is nevertheless environment-specific: it relies on the expected \texttt{kmalloc-1024} layout, exposed kernel symbols and disabled freelist hardening.

\subsection{Comparison with Previous eBPF Exploits}

The most relevant comparison for CVE-2024-58100 is CVE-2017-16995. The root cause is different, but the exploitation strategy is structurally close: both exploits start from a verifier-approved BPF program, obtain a kernel read/write capability, and then use that primitive to corrupt a security-sensitive kernel object.

In CVE-2017-16995, the exploit builds a BPF-based read/write oracle exposed to userspace through a command channel. The attacker can request operations such as leaking pointers, reading arbitrary 64-bit values, or writing arbitrary 64-bit values. The exploit then uses networking-related kernel objects to locate credentials: it leaks a \texttt{sk\_buff}, follows it to the associated \texttt{struct sock}, searches for the \texttt{sk\_rcvtimeo} sentinel, and recovers \texttt{sk\_peer\_cred}. Once a valid \texttt{struct cred} is found, the exploit overwrites UID/GID fields and grants capabilities, finally returning to userspace with root privileges.

CVE-2024-58100 follows a comparable post-exploitation logic, but with a different memory primitive and final target. Instead of forging an arbitrary read/write primitive and directly patching \texttt{struct cred}, it obtains a stale packet pointer through the global-subprogram compatibility bug. The resulting UAF write poisons the \texttt{kmalloc-1024} SLUB freepointer and forces a later pipe-buffer allocation into a controlled userspace page. The exploit then replaces the first buffer's \texttt{ops} field, triggers the fake release path and modifies \texttt{modprobe\_path}.

The main similarity is therefore not the vulnerable verifier function, but the exploit architecture: both attacks convert a verifier/runtime mismatch into a reusable kernel memory primitive and then use userspace orchestration to drive the final escalation stage. CVE-2017-16995 uses the primitive to locate and patch credentials directly; CVE-2024-58100 manipulates allocator metadata and redirects kernel execution indirectly through a forged \texttt{pipe\_buffer} object.

\subsection{Patch and Mitigation}

The purpose of the fix is to make EXT-program attachment respect the \texttt{changes\_pkt\_data} property of the global subprogram being replaced. After the patch, the verifier computes and stores the \texttt{changes\_pkt\_data} property before validating the \texttt{EXT} attachment. If the replacement program and the replaced global subprogram are incompatible with respect to packet-data modification, the attachment is rejected. This prevents a caller from being verified against a benign packet-data contract and then executing a replacement implementation with stronger packet-mutating side effects.

Conceptually, the patched verifier behaves as follows:

\begin{lstlisting}[caption={Conceptual model of the fix.}]
ret = check_cfg(env);
if (ret < 0)
    goto skip_full_check;

ret = check_attach_btf_id(env);
if (ret)
    goto skip_full_check;

if (extension_aux->changes_pkt_data != target_aux->changes_pkt_data)
    return -EINVAL;
\end{lstlisting}

After the fix, the verifier computes \texttt{changes\_pkt\_data} before validating the EXT-program attachment. Therefore, an extension program whose packet-data side effects differ from the replaced global subprogram is rejected. This prevents the caller from being verified under assumptions that no longer hold at runtime.

The practical mitigations are:

\begin{itemize}
    \item update to Linux 6.6.90, 6.12.25, 6.13 or another kernel containing the fix;
    \item restrict \texttt{CAP\_BPF}, \texttt{CAP\_NET\_ADMIN} and the ability to attach EXT programs;
    \item keep \texttt{kptr\_restrict} enabled and avoid exposing kernel symbol addresses to untrusted processes;
    \item disable unprivileged BPF where it is not required.
\end{itemize}

\subsection{Security Impact}

CVE-2024-58100 shows that eBPF verifier correctness does not only depend on local instruction validation. It also depends on preserving side-effect compatibility across function boundaries and dynamic replacement boundaries. Global subprograms make BPF development modular, but they also introduce a contract-based verification model. If an \texttt{EXT} program is allowed to replace a global subprogram while changing the packet-data invalidation contract, the verifier can approve a caller whose runtime packet pointer becomes stale.

The key lesson is that stale verifier state can be as dangerous as a direct bounds-check failure. Here, the stale state is not a scalar range or a stack slot, but a packet pointer whose lifetime ended inside an opaque global subprogram. That stale \texttt{PTR\_TO\_PACKET} becomes a real Use-After-Free primitive.

In the documented laboratory environment, the equivalent locally reproduced primitive was converted into full local privilege escalation. By poisoning the SLUB freepointer, the exploit controlled the placement of a \texttt{pipe\_buffer}, redirected its release callback and overwrote \texttt{modprobe\_path}. The official vulnerability explains how an incompatible EXT replacement can produce the stale-pointer condition; the local program demonstrates what can be achieved once that condition is available. The final chain additionally depends on allocator layout and deliberately weakened hardening settings.

\end{document}
